Large language models (LLMs) are increasingly consulted as an accessible, low-cost source of mental-health guidance. This dissertation examines whether their responses to standardised queries vary with the city and country stated by the user. The dataset contains 1{,}120 responses from seven open-weight and commercial LLMs, generated for 20 cities across four World Bank income categories and six WHO regions using eight fixed scenarios.

The study combines transparent rule-based coding with human judgement. The first versions of all three outcomes failed a pre-specified $F_1$ threshold. Human labels exposed a circular localisation component that had survived agreement between two automated rule sets. Actionability and localisation were redesigned on 112 development labels and then tested, frozen, on a fresh blinded 160-response hold-out. Neither passed the pre-specified hold-out criteria: actionability achieved weighted $\kappa=0.462$, while surface localisation showed specificity of 0.038 despite high sensitivity. Religious framing was assessed separately. This validation gate is central to the interpretation of the results.

The frozen measures produce three sample-specific patterns, but their evidential status differs. Actionability rises across income categories, yet component and joint-model sensitivity analyses indicate that this association is largely explained by crisis-contact inclusion and documented service availability; because actionability failed hold-out validation, RQ1 is exploratory. Surface localisation is higher where a documented crisis service exists and the direction persists when only reference-matched contacts receive credit, but localisation also failed validation, so RQ2 is exploratory. Religious or spiritual recommendations vary by WHO region with separate validation support, although the fitted model does not separate region from income. With only 20 geographic units, none of these associations is global or causal evidence.

An unplanned audit found crisis-contact numbers with visibly suspicious patterns, including fictional 555 exchanges and repeated or sequential digit runs. Their observed rate increased from approximately 1\% in the established-service tier to 34\% in the no-documented-service tier. This is a lower-bound pattern flag, not a fabrication estimate: unmatched numbers remain unknown unless a source establishes that they are correct or incorrect.

The study provides evidence that stated location is associated with differences in LLM mental-health guidance within this 20-city sample. Its strongest contribution is methodological: automated measures can agree with one another while failing construct validation, so rule development, hold-out validation and confirmatory testing must be kept separate. The crisis-contact audit identifies a safety-relevant signal that warrants complete source-based verification and replication.
